\scanpage{14}
% archive sheet number: 13
% lines 1–2
\deleted{\unsure{Prop}} \uline{Lemme B} Supposons \deleted{de plus} que $X$ satisfait les conditions suivantes : (i) (ii) \uline{et}

% lines 3
(iv) Les $\underline O_{X,x}$ $(x\in X)$ sont \uline{noethériens}

% lines 4–10
(v) Pour tt \deleted{\unsure{sous-module}} Idéal coh. $J$, \deleted{\unsure{tout point}} posant $Y=(\operatorname{Supp}\underline O_X/J,\underline O_X/J)$, et tt $y\in Y$ tel que $\underline O_{Y,y}$ est \uline{intègre}, il existe un voisinage ouvert $U$ de $y$ \inserted{dans $Y$} ayant la propriété suivante : toute section $f$ de $\underline O_Y$ sur \deleted{un voisinage} $U$ qui est \deleted{telle que} nulle au voisinage de $y$, est nulle sur $U$.

% lines 10–13
Alors pour toute \deleted{suite} $F$ coh. sur $X$, \deleted{\unsure{et}} tte suite \inserted{croissante} $F_i$ de sous-faisceaux coh. de $F$ est stationnaire \ldots

% lines 14–15
\uline{Corollaire} Sous les conditions (i) à (v), $A=\Gamma(X,\underline O_X)$ est un anneau \uline{noethérien}.

% lines 16–18
Ces conditions sont satisfaites si $X$ est un \deleted{espace de Stein} \inserted{\unsure{polyèdre|polyèdres|polycylindre} analytique} (compact) \deleted{\unsure{dont on trouve}} sur $\mathbb C$. \deleted{dont} t\unsure{t} \deleted{\textquotedblleft \unsure{presque|faisceau} \unsure{tout|étalé} fermé\textquotedblright{} est loc connexe} \ldots
% slanted
\margin{préciser ce point}
